\subsection{Task 3.1 -- Kiến trúc hệ thống}

\subsubsection{Tổng quan}

Kiến trúc BrewLite được mô tả theo \textbf{C4 ở mức Container} (Hình~\ref{fig:architecture-c4}); mức Context được mô tả bằng lời. Ở mức Context, hệ thống BrewLite phục vụ hai nhóm người dùng: \textbf{Khách hàng} (đặt đồ uống, thanh toán bằng QR chuyển tiền, theo dõi đơn) và \textbf{Nhân viên} có vai trò \texttt{ADMIN} (đối chiếu sao kê, xác nhận thanh toán, duyệt hủy, pha chế). Không có hệ thống bên ngoài nào được tích hợp: ngân hàng giả lập gồm trang chuyển tiền ở frontend và module \texttt{mock-bank} (sổ sao kê) ở backend. Ở mức Container, hệ thống gồm frontend Next.js, backend NestJS (chia module theo domain) và cơ sở dữ liệu PostgreSQL.

\begin{figure}[htbp]
\centering
\includegraphics[width=\linewidth,height=0.9\textheight,keepaspectratio]{architecture-c4.png}
\caption{Sơ đồ kiến trúc BrewLite (C4 Container)}
\label{fig:architecture-c4}
\end{figure}

\subsubsection{Mô tả thành phần}

\paragraph{Các container.}
\mbox{}\\
{\small\renewcommand{\arraystretch}{1.25}
\begin{tabularx}{\linewidth}{@{}l X X@{}}
\toprule
\textbf{Container} & \textbf{Công nghệ} & \textbf{Trách nhiệm} \\
\midrule
Frontend (\texttt{apps/web}) & Next.js (mobile-first/PWA), TailwindCSS, React Query, Zustand & Giao diện khách hàng và trang quản trị; giữ giỏ hàng bằng Zustand; gọi API và theo dõi trạng thái đơn bằng React Query \\
Backend API (\texttt{apps/api}) & NestJS (REST/JSON) & Xử lý nghiệp vụ, xác thực JWT và phân quyền \texttt{CUSTOMER}/\texttt{ADMIN}, truy cập cơ sở dữ liệu \\
PostgreSQL & PostgreSQL + TypeORM & Lưu trữ toàn bộ dữ liệu; migration tự chạy khi API khởi động \\
\bottomrule
\end{tabularx}}

\paragraph{Frontend Next.js.}

Ứng dụng web mobile-first, hướng tới PWA, dùng TailwindCSS cho giao diện. React Query gọi API và polling trạng thái đơn; Zustand giữ giỏ hàng (lưu lại khi tải lại trang). Trang quản trị cho nhân viên nằm trong cùng ứng dụng và chỉ truy cập được với vai trò \texttt{ADMIN}. Frontend cũng tạo mã QR ở trang \texttt{/pay/:paymentId} (thư viện \texttt{qrcode.react}) và có trang chuyển tiền giả lập \texttt{/mock-bank/:paymentId}.

\paragraph{Backend NestJS chia module theo domain.}

Backend là một ứng dụng NestJS (modular monolith). Mỗi module chứa controller, service, entity và DTO của domain mình; module khác chỉ dùng qua service mà module đó xuất ra.

{\small\renewcommand{\arraystretch}{1.25}
\begin{tabularx}{\linewidth}{@{}l X@{}}
\toprule
\textbf{Module} & \textbf{Trách nhiệm} \\
\midrule
\texttt{auth} & Đăng ký, đăng nhập (mật khẩu băm bằng \texttt{bcrypt}), cấp JWT, phân quyền \texttt{CUSTOMER}/\texttt{ADMIN}; tài khoản nhân viên khởi tạo sẵn từ biến môi trường \\
\texttt{products} & Menu, chi tiết sản phẩm, tùy chọn size và topping, giá \\
\texttt{inventory} & Nguyên liệu và công thức; giữ chỗ khi tạo đơn (optimistic locking), hoàn kho khi hủy hoặc hết hạn thanh toán \\
\texttt{orders} & Tạo đơn từ giỏ hàng (Idempotency-Key), máy trạng thái đơn (\texttt{assertTransition}), tự hủy đơn quá hạn, hủy đơn, lịch sử đơn; điểm thưởng (\texttt{order-benefits.service}) \\
\texttt{payments} & Phiên thanh toán, nhân viên xác nhận hoặc từ chối (khóa dòng), tự hoàn tiền khi từ chối \\
\texttt{mock-bank} & Sổ sao kê của ngân hàng giả lập (\texttt{bank\_transactions}, \texttt{BankLedgerService}): ghi tiền vào khi khách chuyển và ghi hoàn tiền \\
\texttt{vouchers} & Khuyến mãi (voucher): mã công khai và mã ẩn, giới hạn lượt dùng \\
\bottomrule
\end{tabularx}}

\paragraph{PostgreSQL, TypeORM và migration.}

Dữ liệu lưu trong PostgreSQL, truy cập bằng TypeORM. Mọi thay đổi schema đi qua migration và migration tự chạy khi API khởi động, nên môi trường nào cũng có schema giống nhau (ADR-003).

\paragraph{Thanh toán mô phỏng và xác nhận của nhân viên.}

Không có callback từ cổng thanh toán. Frontend tạo mã QR ở trang \texttt{/pay/:paymentId}; QR mở trang chuyển tiền giả lập \texttt{/mock-bank/:paymentId} (không cần đăng nhập trên điện thoại), nơi khách chuyển tiền (số tiền cố định, nội dung = mã đơn). Module \texttt{mock-bank} ở backend là sổ sao kê (\texttt{bank\_transactions}, \texttt{BankLedgerService}) ghi tiền vào và hoàn tiền; phiên thanh toán chuyển \texttt{TRANSFER\_RECEIVED} trong khi đơn vẫn \texttt{PENDING}. Nhân viên xác nhận thì đơn thành \texttt{PAID}; từ chối thì thành \texttt{PAYMENT\_FAILED} và hệ thống tự hoàn tiền (ADR-004). Idempotency-Key chống tạo đơn hoặc phiên thanh toán trùng; optimistic locking chống bán vượt tồn kho; khóa dòng (pessimistic locking) chống hai nhân viên cùng xác nhận hoặc hoàn tiền hai lần (ADR-005).

\paragraph{Logging, health check và cấu hình.}

\begin{itemize}
  \item \textbf{Logging:} A1 chưa ghi log từng yêu cầu; hiện chỉ dùng \texttt{Logger} của NestJS ở hai chỗ (tạo tài khoản admin và tác vụ tự hủy đơn quá hạn). Ở Task 5, log được ghi ra stdout để nền tảng triển khai (Render) thu thập.
  \item \textbf{Health check:} A1 hiện chỉ có healthcheck trong \texttt{docker-compose}; endpoint \texttt{GET /health} của API sẽ được làm ở Task 5.2.
  \item \textbf{Đóng gói và CI/CD:} đóng gói bằng Docker; \texttt{docker-compose.yml} chạy frontend, backend và PostgreSQL ở local; pipeline CI/CD chạy Lint $\rightarrow$ Test $\rightarrow$ Build trước khi merge; môi trường production dùng HTTPS (NFR-04). Image Docker được build ở CI và lưu trên GitHub Container Registry (\texttt{ghcr.io}), triển khai lên Render (ADR-008).
  \item \textbf{Cấu hình:} toàn bộ qua biến môi trường, có mẫu trong \texttt{.env.example}; bí mật (khóa JWT, mật khẩu cơ sở dữ liệu) không nằm trong repo.
\end{itemize}

\subsubsection{Luồng chính qua các thành phần}

\begin{enumerate}
  \item Khách đăng nhập; frontend nhận JWT và gửi kèm ở mọi yêu cầu (\texttt{auth}).
  \item Khách xem menu và chọn tùy chọn (\texttt{products}), giỏ hàng nằm ở Zustand.
  \item Frontend gọi \texttt{POST /orders} kèm Idempotency-Key; \texttt{orders} tạo đơn \texttt{PENDING} và \texttt{inventory} giữ chỗ nguyên liệu trong cùng một giao dịch.
  \item Frontend gọi \texttt{POST /payments} kèm Idempotency-Key để tạo phiên thanh toán, rồi tạo mã QR ở trang \texttt{/pay/:paymentId} (thư viện \texttt{qrcode.react}).
  \item Khách quét QR, mở trang \texttt{/mock-bank/:paymentId} và chuyển tiền; \texttt{mock-bank} (\texttt{BankLedgerService}) ghi tiền vào sao kê, phiên thanh toán thành \texttt{TRANSFER\_RECEIVED}.
  \item Nhân viên (\texttt{ADMIN}) xác nhận hoặc từ chối; đơn thành \texttt{PAID} hoặc \texttt{PAYMENT\_FAILED} (kèm hoàn tiền).
  \item Frontend tự làm mới (polling) trạng thái đơn cho khách; nhân viên chuyển đơn \texttt{PAID} $\rightarrow$ \texttt{PREPARING} $\rightarrow$ \texttt{READY} $\rightarrow$ \texttt{COMPLETED}.
\end{enumerate}

\subsubsection{Lý do chọn kiến trúc, ưu điểm, hạn chế và trade-off}

\paragraph{Lý do chọn.}

BrewLite phục vụ một cửa hàng duy nhất, do nhóm sinh viên làm trong một học kỳ. Một \textbf{modular monolith} trên một monorepo cho ranh giới nghiệp vụ rõ, chia việc theo module được, giao dịch cơ sở dữ liệu đơn giản và triển khai chỉ cần Docker, CI/CD và HTTPS. Kiến trúc dịch vụ vi mô không cần thiết ở quy mô này.

\paragraph{Ưu điểm.}

\begin{itemize}
  \item Dễ chia việc và review: mỗi người phụ trách một hoặc vài module, ít xung đột mã.
  \item Giao dịch nhiều bước (đặt đơn và giữ chỗ kho, hoàn tiền) làm được trong một giao dịch cơ sở dữ liệu.
  \item Triển khai đơn giản: một API, một frontend, một cơ sở dữ liệu.
  \item Ranh giới module rõ nên sau này có thể tách dịch vụ nếu cần.
  \item Tài liệu và code cùng một repo, dễ truy vết.
\end{itemize}

\paragraph{Hạn chế.}

\begin{itemize}
  \item Các module chạy chung tiến trình: lỗi nặng ở một module có thể ảnh hưởng toàn hệ thống, không scale riêng từng module.
  \item Thanh toán chưa tự động: nhân viên phải đối chiếu sao kê và xác nhận thủ công.
  \item Polling (khoảng 4 giây ở trang quản trị) làm cập nhật trạng thái có độ trễ và tạo tải thừa khi nhiều người mở màn hình theo dõi.
  \item JWT khó thu hồi trước hạn; migration tự chạy khi khởi động không phù hợp nếu có nhiều bản API cùng khởi động.
\end{itemize}

\paragraph{Trade-off.}
\mbox{}\\
{\small\renewcommand{\arraystretch}{1.25}
\begin{tabularx}{\linewidth}{@{}l X X l@{}}
\toprule
\textbf{Quyết định} & \textbf{Được} & \textbf{Mất} & \textbf{ADR} \\
\midrule
Monorepo & Thay đổi cắt ngang nằm trong một Pull Request; tài liệu cùng phiên bản với code & CI chạy rộng, không phân quyền riêng từng phần & ADR-001 \\
Module theo domain & Chia việc rõ, giao dịch đơn giản, triển khai một tiến trình & Lỗi lan chung, không scale riêng & ADR-002 \\
PostgreSQL, TypeORM, migration tự chạy & Schema có phiên bản, dựng môi trường nhanh & Rủi ro khi nhiều bản API khởi động cùng lúc & ADR-003 \\
Ngân hàng giả lập, nhân viên xác nhận & Đúng yêu cầu giảng viên, không phụ thuộc dịch vụ ngoài & Đối chiếu thủ công, khác cổng thật & ADR-004 \\
Idempotency-Key, optimistic và pessimistic locking & Không đơn hoặc thanh toán trùng, không bán vượt tồn kho, không hoàn tiền hai lần & Frontend giữ khóa; khóa lưu vĩnh viễn; khóa dòng làm yêu cầu đến sau phải chờ & ADR-005 \\
JWT, \texttt{CUSTOMER}/\texttt{ADMIN} & Không lưu phiên phía server, phân quyền tập trung & Khó thu hồi token & ADR-006 \\
Polling thay WebSocket & Đơn giản, ít hạ tầng & Trễ cập nhật, tải thừa & ADR-007 \\
\bottomrule
\end{tabularx}}

\subsubsection{Danh mục quyết định kiến trúc (ADR)}

Các ADR nằm trong thư mục \texttt{docs/architecture/adr/}.

{\small\renewcommand{\arraystretch}{1.25}
\begin{tabularx}{\linewidth}{@{}l X X@{}}
\toprule
\textbf{Mã} & \textbf{Quyết định} & \textbf{Tên tệp} \\
\midrule
ADR-001 & Dùng một repo (monorepo) cho web, API và tài liệu & 001-monorepo.md \\
ADR-002 & Backend NestJS chia module theo domain & 002-nestjs-modules-theo-domain.md \\
ADR-003 & PostgreSQL, TypeORM và migration tự chạy & 003-postgresql-typeorm-migration.md \\
ADR-004 & Ngân hàng giả lập và nhân viên xác nhận thanh toán & 004-ngan-hang-gia-lap-nhan-vien-xac-nhan.md \\
ADR-005 & Idempotency-Key và cơ chế khóa (optimistic cho tồn kho, pessimistic cho xác nhận thanh toán) & 005-idempotency-key-va-optimistic-locking.md \\
ADR-006 & JWT và phân quyền \texttt{CUSTOMER}/\texttt{ADMIN} & 006-jwt-va-phan-quyen-customer-admin.md \\
ADR-007 & Polling thay WebSocket & 007-polling-thay-websocket.md \\
\bottomrule
\end{tabularx}}

Số ADR-008 dành cho quyết định chọn nền tảng triển khai ở Task 3.2.
